\subsection{Different Notations for the Different Frameworks}
\label{sec:frameworks:notation}

\begin{table}[]
    \centering
    \begin{tblr}{colspec = {*{6}{Q[l,m]}},
                 row{1} = {font=\bfseries, bg=gray!50},
                 row{odd[2]} = {bg=gray!25},
                 rowsep=4pt
                 }
        \toprule
        CUDA             & HIP              & OpenCL            & SYCL              & Kokkos                    & OpenMP   \\\midrule
        thread           & thread           & work-item         & work-item         & thread                    & thread   \\
        warp             & wavefront        & sub-group         & sub-group         & --                        & --       \\
        {thread\\block}  & {thread\\block}  & work-group        & work-group        & team                      & team     \\
        grid             & grid             & nd-range          & nd-range          & league                    & league   \\\midrule
        register         & register         & {private\\memory} & {private\\memory} & register                  & register \\
        {shared\\memory} & {shared\\memory} & {local\\memory}   & {local\\memory}   & {scratchpad\\memory}      & shared   \\
        {global\\memory} & {global\\memory} & {global\\memory}  & {global\\memory}  & \makecell{memory\\spaces} & --       \\
        \bottomrule
    \end{tblr}
    \caption{%
    The different notations for the different portability languages and frameworks. 
    Since the greatest focus of this dissertation lies on SYCL, I will use its naming conventions for the remainder of this dissertation. 
    For Standard \cpp, these notations do not apply since the \cpp standard has no notion of the underlying hardware because it is only defined on an abstract state machine. 
    For \gls{hpx}, the notation depends on the used executor. 
    }
    \label{tab:frameworks_notations}
\end{table}

In the previous sections, I introduced the programming frameworks I will use in this dissertation. 
\autoref{tab:frameworks_notations} shows the different notations between the programming frameworks. 
Since \gls{hip} is extremely closely related to \gls{cuda}, both programming frameworks use the same notations. 
Originally, SYCL was based on \gls{opencl} and, therefore, also inherited its notation. 
On the other hand, Kokkos uses again another notation more closely related to the \gls{openmp} target offloading notation. 

Since it is not feasible to always mention all notations and the main programming framework used in this dissertation is SYCL, from now on, I will stick to the SYCL/\gls{opencl} notations. 
If you are more familiar with the \gls{cuda} or Kokkos notations, \autoref{tab:frameworks_notations} can be used as a translation aid. 

All these notations are directly based on the hardware structure of modern \glspl{gpu} and form a hierarchy. 
A \textit{work-item} is the smallest unit of execution. 
However, on modern \glspl{gpu}, a number of these units are organized in small groups called a \textit{sub-group} that are executed in lock-step. 
Several such sub-groups are further grouped into a \textit{work-group}. 
All work-items in a work-group can share the same memory and can be synchronized independently of work-items in other work-groups. 
The whole execution range of an application is represented by an \textit{nd-range}.

A similar hierarchy can be seen when looking at modern \gls{gpu} memory types. 
The \textit{private memory} is the fastest memory but only has a rather limited capacity. 
On the other hand, \textit{private memory} is a bit slower but larger. 
All work-items in a work-group share the same private memory which is also sometimes called \enquote{manually managed cache}. 
The \textit{global memory} is the actual device memory marketed by \gls{gpu} vendors and compared to the other memory types quiet slow. 
